\documentclass[12pt,a4paper]{article}

\usepackage[utf8]{inputenc}
\usepackage{amsmath}
\usepackage{graphicx}
\usepackage{hyperref}
\usepackage{listings}
\usepackage{color}
\usepackage{geometry}

\geometry{a4paper, margin=1in}

\definecolor{lightgray}{rgb}{0.95,0.95,0.95}
\lstset{
    backgroundcolor=\color{lightgray},
    basicstyle=\ttfamily\small,
    breaklines=true,
    frame=single
}

\title{\textbf{AudioSpread: Real-Time Audio-to-Sheet-Music Architecture}}
\author{Project Documentation}
\date{\today}

\begin{document}

\maketitle
\tableofcontents
\newpage

\section{Introduction}
The \textbf{AudioSpread} project is an end-to-end machine learning architecture designed to transcribe raw acoustic audio signals (e.g., live ukulele or piano music) into standard Western sheet music notation in real-time. The architecture is explicitly optimized for consumer-grade hardware, ensuring millisecond latency without relying on heavy cloud compute. 

\section{Architecture Overview}
The system is bifurcated into three distinct modules: the Real-Time Inference Backend, the Web-based Visualizer Frontend, and an optional Training Pipeline.

\subsection{Inference Backend (\texttt{server.py})}
The backend relies on \textbf{FastAPI} and \textbf{NumPy/Librosa}. It continuously receives rolling 4096-sample buffers over a WebSocket endpoint (\texttt{/ws}). Pitch extraction relies on a custom \texttt{NoteTracker} class utilizing a Fast Fourier Transform (FFT) combined with parabolic interpolation. 

Crucially, the tracker operates on a strict three-state machine (\texttt{SILENCE $\rightarrow$ ATTACK $\rightarrow$ LOCKED}). This entirely prevents mid-decay harmonic jumps or octave errors. A confirmed MIDI note is broadcast back to the frontend alongside the detected RMS velocity.

\subsection{Frontend Visualizer (\texttt{transcriber-ui})}
A lightweight \textbf{React.js} client manages microphone capture inside the browser context via the modern \textbf{AudioWorkletNode} API (bypassing native OS library conflicts). The client streams binary `Float32Array` audio to the WebSocket server. 

Incoming JSON MIDI tokens are parsed, and an echo-suppression algorithm ensures consecutive identical pitches aren't redundantly re-triggered. The notes are finally fed into the \textbf{VexFlow} rendering engine, drawing them onto an infinite horizontal scrolling HTML5 Canvas stave. The UI is built with a premium ``glassmorphism'' CSS aesthetic.

\subsection{Training Pipeline (\texttt{train.py})}
An optional training pipeline leverages a Convolutional Recurrent Neural Network (CRNN) to map audio frequencies to MIDI pitches.
\begin{itemize}
    \item \textbf{Dataset:} The pipeline utilizes the MAESTRO v3 dataset.
    \item \textbf{Audio Representation:} Raw waveform audio is transformed into Constant-Q Transforms (CQT).
    \item \textbf{Dynamic Loading:} To prevent RAM exhaustion (OOM errors), the script employs an on-the-fly random cropping strategy and interval-based MIDI caching.
    \item \textbf{Export:} The model is exported to an ONNX graph (\texttt{transcriber\_quantized.onnx}) for inference.
\end{itemize}

\section{Setup and Installation}
\subsection{System Prerequisites}
Audio capture now operates directly within the browser, meaning the system runs natively on Windows PowerShell, macOS, and Linux without requiring WSL or PulseAudio configurations.

\subsection{Environment Bootstrapping}
A Python 3 virtual environment must be isolated for the backend processes:
\begin{lstlisting}[language=bash]
# Windows PowerShell Execution
.\setup.ps1
\end{lstlisting}
The \texttt{setup.ps1} script automatically generates the virtual environment and installs \texttt{torch}, \texttt{librosa}, \texttt{fastapi}, \texttt{uvicorn}, \texttt{websockets}, and \texttt{onnxruntime}.

\section{Features and Capabilities}

\subsection{Strict State-Machine Pitch Detection}
The algorithm enforces a hard rule where a new pitch can only be targeted if the audio signal previously returned to a silent threshold. This eliminates ``messy'' false notes generated during string decay harmonics. 

\subsection{Interactive CLI Training}
The \texttt{train.py} script features a comprehensive CLI via \texttt{argparse}. Omitting arguments spawns an interactive menu that dynamically scans the MAESTRO dataset directories.

\section{Execution Protocol}
The application requires simultaneous execution of the backend WebSocket provider and the React consumer:
\begin{enumerate}
    \item \textbf{Backend:} Execute \texttt{uvicorn server:app --port 8000 --reload}.
    \item \textbf{Frontend:} Execute \texttt{npm start} inside the \texttt{transcriber-ui} directory.
\end{enumerate}

\end{document}
